\section{UNITARY REPRESENTATIONS}

In this section, the vector spaces are over K = R or C.

Recall that a linear representation of a group G in a K-vector space E is a homomorphism \(\varrho\) from G into the group \(\mathbf{GL}(\mathrm{E})\) of automorphisms of E (A, VIII, p. 387, def. 1). One says that E is the space of the representation \(\varrho\) , and that the dimension of E is the dimension of \(\varrho\) , also denoted \(\dim(\varrho)\) .

One can identify a representation of G in a K-vector space with a K{[}G{]}-module (loc. cit.), and one will use the corresponding terminology, for example concerning direct sums of representations, morphisms of representations or the representation on C obtained from a representation in an R-vector space by extension of scalars.

A representation \(\varrho\colon\mathbf{G}\to\mathbf{GL}(\mathbf{E})\) is said to be faithful if it is an injective homomorphism.

\pandocbounded{\includegraphics[keepaspectratio]{images/part003_3d1522a8.jpg}}

Note that this does not mean that the K{[}G{]}-module associated with \(\varrho\) is a faithful K-module (A, II, p. 28).

The character of a representation of G in a finite-dimensional K-vector space is the map \(\chi_{\varrho}\) from G to K such that \(\chi_{\varrho}(g) = \operatorname{Tr}(\varrho(g))\) for all \(g \in G\) , cf. A, VIII, p. 388.

